\section{Problem context}
The gap between human and robots still remains significant and leads to communication issues in robotic environments, the interpretation of commands may well be too rigid and may not translate as intended to the real scenario. Furthermore, the task of coordinating multiple agents in an environment arises more issues including the previous, as the difficulty of synchronizing more than one agent raises exponentially due to integration and coordination issues. In order to tackle that gap, several techniques have been developed, including simulation environments to reduce the \textit{sim2real} gap of robotics deployment. However, there are still some performance issues these frameworks cannot cover. Reason why different artificial intelligence methodologies have started to been applied for robotics. In parallel, the rise of multi modal models including the vision language models (VLMs) has opened a new possibility frontier for bridging this gap. VLMs offer the ability of perceiving visual scenes alongside natural language reasoning simultaneously, extending a more flexible interface for human-robot interaction. Nevertheless, the successful integration of this technology to robot physical execution of multiple tasks remains an open challenge to be explored throughout this project. 

\section{Motivation}
Vision language models had little to no other integration in other field rather than \textit{chatbot} plugins and image captioning. Reason why I have wanted to experiment VLM integration in another such an interesting field as robotics. I will also experimenting with a high-level reasoning engine able to think for itself instead of relying on hard-coded actions and positions to give to the agents. Likewise, this project aims to use state of the art models and frameworks for the experimentation part. 

\section{Objectives}
\subsection{Main objective}
Experimentation with the integration of reasoning capacities of a vision language model in a controlled simulated multi-agent robotic environment to solve cooperative manipulation tasks while covering the whole process alongside the SOTA of the models in the documentation.

\subsection{Specific objectives} % ???
The project also aims to meet certain objectives which are more specific. Among those there are the design of progressively more complex environments in the robotic scenario simulator. Starting from the most basic tasks, advancing to more complex training alternatives which require more sophisticated techniques and sharper model performance. Secondly, another goal to achieve is to implement a communication layer between VLM and the agents, as the vision language model will be hosted on a server while all the agent manipulation will be run locally. Moreover, I will also define metrics to measure the success rate of cooperative tasks and the performance of the VLM. Such as the model's ability to generalize across different environmental variations. 

\section{Project scope}
I will distinct between what the project will include and what will not. Let's go over the topics and ideas the project will execute first. Firstly, the vision language models will be integrated as high-level decision making orchestrator brains. Furthermore, the orchestration pipeline for communication will be custom made. Finally, the practical area of the project is strictly developed inside NVIDIA's Isaac Sim, taking advantage of its complex physics engine and realistic rendering capacities (RTX).

On the other hand, there are several aspects the project will not cover, and I belief it is essential to mention them in order to not create any sort of confusion. The physical display of the robots won't be covered, they will be strictly displayed in the controlled environment. This project will also not train foundational VLMs from scratch, the use models are pre-trained open-source alternatives. 